\centerline{\bf \Huge Комбинаторика}
\title{Комбинаторика}

\begin{flushright}
    \scriptsize Слово <<списать>> и слово <<бан>> для вас означает одно и то же!\\Доктор Анджур. Остров ЛМШ.
\end{flushright}

Комбинаторика $-$ огромнейшая часть Олимпиадного мира. Это такая область математики, которая изучает и решает задачи, связанные с размещением и сочетанием различных предметов. На каждой серьезной Олимпиаде обязательно будет 1-2 задачи по комбинаторике, поэтому если не изучить хорошо этот раздел, то баллы неминуемо потеряются.

\textbf{\Large Четыре всадника комбинаторики}

Базовая часть комбинаторики состоит из 4 правил, отвечающие на 4 вопроса: <<Когда складывать?>>, <<Когда умножать?>>, <<Когда возводить в степень?>> и <<Когда делить?>>. Давайте разберем 4 задачи, в которых разберем, как работают эти правила.\\

\textit{Правило сложения.} Если элемент A можно выбрать n способами, а элемент B можно выбрать m способами, то выбрать A \textbf{или} B можно n + m способами. В этом правиле ключевое слово \textbf{или}. Когда вы хотите вы выбрать <<это или это>>, то количество вариантов складывается.\\


\problem{1}У Расула Владимировича дома 6 кошек $-$ 3 породистые и 3 дворняжки. Сколькими способами Анджур Аркадьевич может выбрать, какую кошку ему погладить?

\textbf{Ответ: 6.}

\textbf{Решение:} Заметим, что Анджуру Аркадьевичу нужно выбрать ИЛИ породистую ИЛИ дворняжку. Количество способов выбрать породистую, конечно, 3. То же самое с дворняжками. По правилу сложения эти способы нужно сложить! Получим ответ 6, который и так был очевиден.\\

\textit{Правило умножения.} Если элемент A можно выбрать n способами, а элемент B можно выбрать m способами, то выбрать A \textbf{и} B можно $n\cdot m$ способами. В этом же правиле ключево слово \textbf{и}. Если вы хотите выбрать один предмет \textbf{и} второй предмет, то количество вариантов выбрать первый нужно умножить на количество вариантов выбрать второй.
\newpage

\problem{2} Анджуру Аркадьевчу стало мало одной кошки, теперь он хочет погладить сразу две кошки. Чтобы было справедливо, он хочет выбрать одну породистую и одну дворняжку. Сколькими способами он сможет это сделать?

\textbf{Ответ: 9.}

\textbf{Решение:} Давайте сначала выберем породистую кошку. Количество способов это сделать ровно 3. Теперь для каждого выбора породистой кошки у нас есть ещё по 3 способа выбрать какую-то дворняжку. Иначе говоря, выбор породистой кошки никак не зависит от выбора дворняжки. Поэтому количество способов будет $3\cdot 3=9$. 

\problem{3}\textit{Когда делить?.} Решим ту же самую задачу, только теперь без условия того, что нужно погладить одну и одну дворняжку. Нужно просто выбрать две кошки. Зовут их Буся, Люся, Дуся, Маруся, Кристи и Ванга.

\textbf{Ответ: 15.}

\textbf{Решение:} Теперь нам нужно выбрать две кошки из 6. Давайте сперва выберем одну, а к ней вторую. Сколькими способами можно выбрать первую кошку? Конечно же, шестью. Осталось 5 кошек. Сколькими способами можно выбрать вторую кошку к первой? Конечно, пятью. Тут работает правило умножения, так как нужно выбрать первую \textbf{И} вторую кошку. Стало быть ответ будет $5 \cdot 6 = 30$. Но всё не так просто, как кажется. Здесь мы посчитали некоторые варианты несколько раз. Нет ведь никакой разницы, какую кошку брать первой, а какую второй, порядок здесь не важен. Если мы возьмем сначала Бусю, а потом Люсю или сначала Люсю, а потом Бусю, то это будет одно и то же. Поэтому количество вариантов нужно разделить на 2. Ответ будет 15.

\problem{4}\textit{Когда возводить в степень?} Анджур Аркадьеьвич на протяжении всей недели приходи к Расулу Владимировичу и гладил одну кошку. Сколькими способами он мог это сделать? 

\textbf{Ответ: $6^7$}

\textbf{Решение:} Каждый день у Анджура Аркадьевича ровно 6 вариантов, какую кошку погладить. Но так как он должен погладить и в понедельник, и во вторник, и в среду и т.д., во все дни недели, то здесь срабатывает правило умножения. Поэтому 6 нужно уможнить 7 раз на себя. То есть ответ $6^7$.
\newpage

\textbf{\Large Цербер комбинаторики}\\

\textbf{\large Перестановки}

\problem{1}Сколькими способами можно выстроить в ряд 6 кошек Расула Владимировича?

\textbf{Ответ: 6!}

\textbf{Решение:} Давайте сначала посмотрим, сколькими способами можно поставить кошку на первое место. Так как все кошки на данный момент свободны, можно поставить любую из них, то есть 6 вариантов. На второе место уже 5 вариантов, так как одна кошка уже занята. На третье место 4 варианта, так как две кошки уже заняты и так далее, на последнее место останется один единственный вариант. Все эти числа нужно умножить, так как работает правило умножения. Получим ответ $6\cdot 5 \cdot 4 \cdot 3 \cdot 2 \cdot 1 = 720$.

На самом-то деле число $n!$ это количество перестановок $n$ предметов. Обозначается как $P_n = n!$

\textbf{\large Размещения}

\problem{2}Гэндо Икари выстроил в ряд 10 потенциальных пилотов и хочет выбрать, кто из них будет первое дитя, кто второе, кто третье, кто четвертое и кто пятое. Сколькими способами он может это сделать?

\textbf{Ответ:$\dfrac{10!}{5!}$}

\textbf{Решение:} Выберем сначала первое дитя. Это можно сделать одним из 10 способов. Далее выберем второе дитя, это можно сделать девятью способами, так как один пилот уже занят. И так далее пятое дитя можно выбрать шестью способами. Получим ответ $10 \cdot 9 \cdot 8 \cdot 7 \cdot 6$. Заметим, что тут не нужно делить на $5!$, так как порядок нам важен.

Выбор k предметов из n предметов с учетом порядка называется размещением. Число размещений обозначается как $A_n^k = \dfrac{n!}{k!}$ 

\textbf{\large Сочетания}

\problem{3} Комнада ЛмШ из 8 человек(Эльзята, Алина, Эрдем, Баир, Юля, Эвена и Амуля) собралась, чтобы найти и победить НеНачальника. Было принято решение, что им нужно разделиться $-$ двое будут искать в Троицком, а шестеро в Бага-Буруле. Сколькими способами они могут разделиться?

\textbf{Ответ:$C_6^2$}

\textbf{Решение:}На самом деле главное выбрать тех, кто поедет в Троицкое, а остальные точно уже поедут в Бага-Бурул. Сколькими способами можно выбрать двух людей 
\newpage
из восьми, которые поедут в Троицкое? Это ведь та же самая задача, что и с шестью кошками, где нужно было выбрать две из них! Ответ будет $8\cdot 7 :2  =28$. 

Выбор $k$ предметов из n предметов \textbf{без учета порядка} называется сочетанием. Число сочетаний обозначается как $C_n^k = \dfrac{n!}{k!(n-k)!}$

\textbf{\Large Анаграммы}

\problem{1}Сколько различных слов можно составить из данных, переставляя в них буквы:\\
а) БУЕРАК\\
б) ДУРНОЙВКУС\\
в) ЭЛЕКТРОФОРЕЗ\\

\textbf{\Large Метод шариков и перегородок}

\problem{1}Кишибо Рохан хочет переплести свои книги 15 книг в красный, синий и зеленый цвета. Сколькими способами он может это сделать?

\problem{2}После неудачных поисков в Троицком и Бага-Буруле команда ЛмШ из восьми человек собралась, чтобы поискать НеНачальника в Элисте, но тот подкрался к ним сзади! Применив способность Tricky Fingers своего стенда, он украл из кошельков ребят все монетки, после чего раскидал по 8 кошелькам 40 фальшивых монет. Сколькими способами он мог это сделать, если\\
а) Ни один кошелек не должен был остаться пустым\\
б) Какие-то кошельки могли остаться пустыми

\textbf{\Large Где моя лега?}

\problem{1}Герман скачал новую игру себе на телефон, в которой есть 300 различных предметов, среди которых 10 легендарных. Так же в игре есть наборы, которые содержат 5 предметов и не более одного легендарного. Герман открыл уже кучу наборов, но все никак не может получить легендарный предмет. Почему?

\textbf{Решение:} Давайте посчитаем, с какой вероятностью выпадает легендарный предмет в этом наборе. Чтобы считать вероятность, мы воспользуемся формулой\\
\\
$\dfrac{\text{Количество нужных исходов}}{\text{Количество всех исходов}}\cdot 100\%$. Давайте посчитаем количество всех исходов. Это количество вариантов взять 5 предметов из 3000 предметов, то есть $C_{3000}^5$. Теперь посчитаем количество нужных исходов. Нужные исходы для нас это когда нам попался легендарный предмет. Пусть он, действительно, попался и осталось добавить в набор ещё 4 нелегендарных предмета из оставшихся 2990 предметов, то есть $C_{2990}^4$. Осталось подставить в формулу и получим, что шанс выпадения легендарного предмета $\dfrac{C_{2990}^4}{C_{3000}^5}\cdot100\% = 0.16\%$. Неудивительно, что Герман не смог ничего выбить.